\documentclass[12pt]{article}
\usepackage{amsmath,amssymb,amsfonts}
\usepackage[margin=1in]{geometry}

\begin{document}

\section*{Problem 5.7}

\textbf{Problem:}
\begin{itemize}
\item[(a)] Show that the energy transferred to a spring by a time-dependent force $F(t)$ defined by (Fig.\ 5.6):
\[
F(t) = F_0(2t/\tau + 1/2), \quad |t| \leq \tau, \quad \text{is} \quad \Delta E = \frac{1}{2}F_0^2k^{-1}\left(\frac{\sin\omega\tau}{\omega\tau}\right)^2.
\]
\emph{Hint:} Take derivatives of $F(t)$ until you reach a function whose Fourier transform is known.

\item[(b)] Show that the energy transfer in the adiabatic and sudden limits agree with those determined for the probability integral force function (Example 3, Section 5.7).

\item[(c)] For certain finite values of $\omega\tau$ there will be no energy transfer to the spring. Determine these values and explain their origin physically.
\end{itemize}

\bigskip
\textbf{Solution:}

\subsection*{Part (a)}

The energy transferred to a harmonic oscillator of spring constant $k$ and natural frequency $\omega = \sqrt{k/m}$ by a force $F(t)$ is given by (Eq.\ 5.6 of the text):
\[
\Delta E = \frac{|\tilde{F}(\omega)|^2}{2m},
\]
where $\tilde{F}(\omega)$ is the Fourier transform of $F(t)$.

The force function is a linear ramp:
\[
F(t) = \begin{cases} 0 & t < -\tau, \\ F_0\left(\frac{t}{2\tau} + \frac{1}{2}\right) & -\tau \leq t \leq \tau, \\ F_0 & t > \tau. \end{cases}
\]

Taking the first derivative:
\[
F'(t) = \begin{cases} F_0/(2\tau) & -\tau < t < \tau, \\ 0 & |t| > \tau. \end{cases}
\]
with step discontinuities: a jump of $+F_0/2$ at $t=-\tau$ (not relevant for derivative) --- actually, let us be more careful. We have:
\[
F'(t) = \frac{F_0}{2}\delta(t+\tau) + \frac{F_0}{2\tau}[\Theta(t+\tau)-\Theta(t-\tau)] + \frac{F_0}{2}\delta(t-\tau),
\]
but it is simpler to use the second derivative approach.

Taking the second derivative:
\[
F''(t) = \frac{F_0}{2}\delta(t+\tau) - \frac{F_0}{2}\delta(t-\tau) + \frac{F_0}{2\tau}[\delta(t+\tau) - \delta(t-\tau)].
\]

Actually, let us use a cleaner approach. Define $F(t)$ as the ramp from $0$ to $F_0$ over $[-\tau, \tau]$. The Fourier transform of the energy formula gives:
\[
\tilde{F}(\omega) = \frac{1}{\sqrt{2\pi}}\int_{-\infty}^{\infty}F(t)e^{i\omega t}\,dt.
\]

Since $F(t)$ is not square-integrable (it approaches $F_0$ for $t > \tau$), we work with the impulse. The key quantity is:
\[
\Delta E = \frac{1}{2m}\left|\int_{-\infty}^{\infty}F(t)e^{i\omega t}\,dt\right|^2.
\]

For the ramp force, using the derivative method: if $F'(t) = \frac{F_0}{2\tau}$ for $|t|<\tau$ and zero otherwise (this is a rectangular pulse), then:
\[
\widetilde{F'}(\omega) = \frac{F_0}{2\tau}\int_{-\tau}^{\tau}e^{i\omega t}\,dt = \frac{F_0}{2\tau}\cdot\frac{2\sin\omega\tau}{\omega} = \frac{F_0\sin\omega\tau}{\omega\tau}.
\]

Since $\widetilde{F'}(\omega) = -i\omega\tilde{F}(\omega)$ (up to the static part), we get:
\[
|\tilde{F}(\omega)|^2 = \frac{|\widetilde{F'}(\omega)|^2}{\omega^2} = \frac{F_0^2\sin^2\omega\tau}{\omega^4\tau^2}.
\]

The energy transferred is:
\[
\Delta E = \frac{|\tilde{F}(\omega)|^2}{2m} = \frac{F_0^2\sin^2\omega\tau}{2m\omega^4\tau^2}.
\]

Using $k = m\omega^2$:
\[
\boxed{\Delta E = \frac{F_0^2}{2k}\left(\frac{\sin\omega\tau}{\omega\tau}\right)^2 = \frac{1}{2}F_0^2 k^{-1}\left(\frac{\sin\omega\tau}{\omega\tau}\right)^2.}
\]

\subsection*{Part (b)}

\textbf{Adiabatic limit} ($\omega\tau \to \infty$, i.e., slow turn-on):
\[
\Delta E \to \frac{F_0^2}{2k}\cdot 0 = 0.
\]
No energy is transferred---the system adjusts quasi-statically, consistent with the adiabatic theorem.

\textbf{Sudden limit} ($\omega\tau \to 0$, i.e., instantaneous turn-on):
\[
\frac{\sin\omega\tau}{\omega\tau} \to 1, \qquad \Delta E \to \frac{F_0^2}{2k}.
\]
This is the maximum energy transfer, equal to the potential energy stored in the spring under static deflection $F_0/k$, which matches the sudden limit result from the probability integral force function.

\subsection*{Part (c)}

The energy transfer vanishes when $\sin\omega\tau = 0$ (with $\omega\tau \neq 0$), i.e., when:
\[
\boxed{\omega\tau = n\pi, \qquad n = 1, 2, 3, \ldots}
\]

\textbf{Physical explanation:} The condition $\omega\tau = n\pi$ means $\tau = nT/2$ where $T = 2\pi/\omega$ is the natural period of the oscillator. During the ramp-up time $2\tau$, the oscillator completes exactly $n$ full half-cycles. The energy delivered during the first half of the ramp is exactly canceled by the energy removed during the second half, resulting in zero net energy transfer. This is a resonance-avoidance condition.

\end{document}
